%%AMC:latex_engine=lualatex
% ↑ AMC 本体にこのファイルを LuaLaTeX で処理させる指定（消さないでください）
%
% AIビジネス入門 第08回 AIプラットフォームビジネスの全体像と課題 確認問題（CC BY-SA 4.0）
% 既存の資料には問題がなかったため、新しく作った（lesson_plan.md の「確認問題の設計」）
\documentclass[a4paper,11pt]{ltjsarticle}
\usepackage{lfs-fonts}
\usepackage{amsmath,amssymb}
\usepackage{needspace}
\usepackage[nowatermark]{automultiplechoice}
\AMCrandomseed{260801}
\def\AMCbeginQuestion#1#2{\par\noindent{\bfseries 問#1}\quad #2}

\begin{document}

\element{aib08}{
  \Needspace{8\baselineskip}
  \begin{question}{AIB08-Q1}
    プラットフォームビジネスの特徴としてもっとも適切なものはどれですか？
    \begin{choices}
      \wrongchoice{メーカーから消費者へ、一方通行で商品を届ける}
      \correctchoice{利用者どうしが価値を交換する「場」を提供する}
      \wrongchoice{在庫を多く持ち、ほかの店より安く商品を売る}
      \wrongchoice{広告を出さず、利用者の会費だけで運営する}
    \end{choices}
  \end{question}
}

\element{aib08}{
  \Needspace{8\baselineskip}
  \begin{question}{AIB08-Q2}
    「ネットワーク効果」の説明としてもっとも適切なものはどれですか？
    \begin{choices}
      \wrongchoice{広告を増やすほど、サービスの売上が増えること}
      \wrongchoice{たくさん作るほど、1つあたりの費用が下がること}
      \correctchoice{利用者が増えるほど、利用者にとっての価値が高まること}
      \wrongchoice{社員が増えるほど、新しい機能を早く開発できること}
    \end{choices}
  \end{question}
}

\element{aib08}{
  \Needspace{8\baselineskip}
  \begin{question}{AIB08-Q3}
    AI によってプラットフォームビジネスはどのように進化しましたか？ もっとも適切なものを選んでください。
    \begin{choices}
      \wrongchoice{すべての利用者に、同じ内容を見せるようになった}
      \wrongchoice{担当者が手作業で、表示の順番を決めるようになった}
      \wrongchoice{一度決めた組み合わせを、変えずに使うようになった}
      \correctchoice{利用者ごとに、見せる内容を変えるようになった}
    \end{choices}
  \end{question}
}

\element{aib08}{
  \Needspace{8\baselineskip}
  \begin{question}{AIB08-Q4}
    MVP（Minimum Viable Product）の説明としてもっとも適切なものはどれですか？
    \begin{choices}
      \correctchoice{核となる価値をもつ、必要最小限の機能の製品}
      \wrongchoice{考えられるすべての機能をそろえた完成品}
      \wrongchoice{競合よりも価格をもっとも安くした製品}
      \wrongchoice{市場でいちばんよく売れている人気の製品}
    \end{choices}
  \end{question}
}

\element{aib08}{
  \Needspace{8\baselineskip}
  \begin{question}{AIB08-Q5}
    AI プラットフォームのプライバシーの対策としてもっとも適切なものはどれですか？
    \begin{choices}
      \wrongchoice{集めたデータの使い道を、利用者に知らせずに使う}
      \correctchoice{利用目的を示し、利用者がオプトアウトできるようにする}
      \wrongchoice{処理を速くするため、データを暗号化せずに保存する}
      \wrongchoice{プライバシーポリシーは、一度決めたら見直さない}
    \end{choices}
  \end{question}
}

\onecopy{1}{
  \noindent{\large\bfseries AIビジネス入門　第08回　確認問題}\hfill
  氏名\,\underline{\hspace{45mm}}\par\medskip
  \AMCcodeGridInt[h]{student}{8}
  \medskip
  \insertgroup{aib08}
}

\end{document}
